\section{有效作用量}

引入双线性源 $J$，使配分函数
\begin{equation}
  Z[J]
  =\mathrm{Tr}\,
  T_\tau\exp\Bigl(
  -\int_0^\beta(H-J\cdot\psi^\dagger\psi)
  \Bigr)
\end{equation}
成为连通格林函数的生成泛函：
\begin{equation}
  W[J]=\ln Z[J],
  \qquad
  G=\frac{\delta W}{\delta J}.
\end{equation}
Legendre 变换定义有效作用量（1PI 生成泛函）
\begin{equation}
  \Gamma[G]
  =W[J]-\mathrm{Tr}(JG),
\end{equation}
满足
\begin{equation}
  \frac{\delta\Gamma}{\delta G}=-J.
\end{equation}
无源物理点 $J=0$ 对应 $\delta\Gamma/\delta G=0$，即自洽确定物理格林函数。对自由理论，
\begin{equation}
  \Gamma_0[G]
  =\mathrm{Tr}\ln(-G^{-1})-\mathrm{Tr}\bigl((G_0^{-1})G\bigr)+\mathrm{const}.
\end{equation}

\begin{note}
有效作用量的二阶导数给出不可约顶点；Dyson、Bethe--Salpeter 等方程均可由 $\Gamma$ 的驻点与展开得到。
\end{note}
